\ChapterImageStar[cap:dar]{Análisis de Decisiones y Resolución}{./images/fondo.png}\label{cap:dar}
\mbox{}\\


\section{Metodología de evaluación}

La metodología de evaluación empleada para la selección de las notaciones y lenguajes de modelado arquitectónico en el proyecto~\NDITI~fue la \textit{Decision Analysis and Resolution} (\DAR) del modelo \CMMI~\citep{CMMIInstitute2010}. Esta metodología permitió analizar de forma estructurada las necesidades del grupo~\GRID~y del entorno institucional, considerando múltiples alternativas de herramientas y enfoques de modelado. El proceso contempló criterios técnicos, metodológicos y pedagógicos, tales como el nivel de estandarización de la notación, la compatibilidad con marcos de arquitectura empresarial, la disponibilidad de herramientas de soporte (por ejemplo, Archi o Enterprise Architect), la curva de aprendizaje, la adopción en contextos académicos y la calidad de la documentación técnica.

\subsection{Lista de criterios}
A continuación se describen los criterios seleccionados para la evaluación.

\subsubsection{Popularidad de la notación}
\textbf{Código:} POPULARIDAD

\textbf{Descripción:} Este criterio evalúa el nivel de adopción y visibilidad de cada notación en la comunidad académica 
y profesional. Para ello, se toma como referencia las tendencias de búsqueda en Google Trends para indicar de manera 
aproximada el interés global que tiene la notación en distintos contextos educativos, industriales y de investigación.

\textbf{Justificación:} En el contexto del proyecto~\NDITI~y del programa de Ingeniería de Sistemas y Computación, 
la popularidad de una notación resulta relevante porque influye en la disponibilidad de ejemplos, tutoriales 
y recursos formativos accesibles para estudiantes y docentes. Una notación con mayor presencia en la comunidad facilita 
el aprendizaje autónomo, la resolución de dudas y aumenta las probabilidades de la vigencia de la notación durante el ciclo 
de vida académico. Si bien esto no garantiza la calidad técnica, representa un valor añadido en términos de soporte
y ecosistema de la notación, por lo que se le considera un criterio pertinente en el proceso de selección.

Para la búsqueda realizada en Google Trends se estableció un rango de tiempo de 2020 a 2024, en línea con el período definido en el SMS.

% TABLA Popularidad según Google Trends.
\begin{table}[H]
	\centering
	\renewcommand{\arraystretch}{1.4} % Espaciado reducido
	\fontsize{9pt}{8pt}\selectfont % Tamaño de fuente 9pt
	\caption{Resultados de búsqueda de popularidad por notación.}
	\label{tab:popularidad_notaciones}
	\begin{tabular}{|>{\centering\arraybackslash}p{2.0cm}|>{\centering\arraybackslash}p{6.5cm}|>{\centering\arraybackslash}p{2.5cm}|>{\centering\arraybackslash}p{2.0cm}|}
		\hline
		{\textbf{Notación}} & {\textbf{Término de búsqueda}}                           & {\textbf{Media de resultados}} & {\textbf{Cuantil [1..3]}} \\
		\hline
		BPMN           & ``BPMN'' 																		& 16                                 & 3                         \\
		\hline
		SysML              & ``SysML''    																& 3                                & 1                         \\
		\hline
		UML                  & ``UML''        															& 64                                & 3                         \\
		\hline
		Archimate           & ``Archimate'' 															& 2                                & 1                        \\
		\hline
	\end{tabular}
	\vspace{5pt}
\end{table}


\subsubsection{Documentación de la notación}
\textbf{Código:} DOC

\textbf{Descripción:} Una documentación sólida es fundamental para garantizar que tanto los estudiantes como los 
docentes puedan aprender y aplicar la notación sin depender exclusivamente de expertos o del conocimiento tácito.
En el ámbito universitario, donde la autonomía en el aprendizaje es un factor determinante, disponer de un material
bien estructurado contribuye a la reducción de errores de diseño e interpretación y contribuye al aumento de calidad 
de los modelos generados. Por esta razón, la documentación constituye un criterio esencial para identificar una notación 
viable en el marco del proyecto~\NDITI.

% TABLA Popularidad según Google Scholar.
\begin{table}[H]
	\centering
	\renewcommand{\arraystretch}{1.4} % Espaciado reducido
	\fontsize{9pt}{8pt}\selectfont % Tamaño de fuente 9pt
	\caption{Resultados de búsqueda de documentación por notación.}
	\label{tab:documentacion_notaciones}
	\begin{tabular}{|>{\centering\arraybackslash}p{2.0cm}|>{\centering\arraybackslash}p{6.5cm}|>{\centering\arraybackslash}p{2.5cm}|>{\centering\arraybackslash}p{2.0cm}|}
		\hline
		{\textbf{Notación}} & {\textbf{Término de búsqueda}}                           & {\textbf{Cantidad de resultados}} & {\textbf{Cuantil [1..3]}} \\
		\hline
		BPMN           & ``BPMN'' 																		& 97100                                 & 2                         \\
		\hline
		SysML              & ``SysML''    																& 31100                                & 2                         \\
		\hline
		UML                  & ``UML''        															& 488000                                & 3                         \\
		\hline
		Archimate           & ``Archimate'' 															& 9260                                & 1                        \\
		\hline
	\end{tabular}
	\vspace{5pt}
\end{table}

\subsubsection{Apalancamiento y formalidad de la notación}
\textbf{Código:} AP FO

\textbf{Descripción:} Este criterio evalúa el grado de formalidad institucional detrás de la creación, 
regulación y mantenimiento de la notación. Se consideran factores como si proviene de un organismo internacional, 
un consorcio académico, una empresa tecnológica o una comunidad abierta, así como la existencia de estándares oficiales, 
certificaciones o procesos de gobernanza.


\textbf{Justificación:} La formalidad con la que una notación es desarrollada influye directamente en su estabilidad, 
coherencia conceptual y continuidad en el tiempo. Para el área de infraestructura de~\TI{}, una notación respaldada por una 
organización sólida brinda mayor confianza en su evolución y reduce el riesgo de quedar obsoleta o fragmentada. 
Adicionalmente una base formal facilita su inclusión en entornos académicos al ofrecer un marco de referencia claro.
Este criterio aporta un nivel importante de seguridad metodológica dentro del proceso de selección de la notación.

\subsubsection{Vigencia de la notación}
\textbf{Código:} VIGENCIA

\textbf{Descripción:} Este criterio analiza la frecuencia y consistencia de las actualizaciones que recibe la notación, 
así como el nivel de actividad de sus desarrolladores o comunidades asociadas. La vigencia se evalúa según la rapidez con 
la que la notación se adapta a nuevas necesidades del sector a través de la actualizacion constante de sus versiones.


\textbf{Justificación:} En un entorno tecnológico en constante cambio como la arquitectura de~\TI{}, es fundamental 
utilizar una notación que no este desactualizada. Para el proyecto~\NDITI, la vigencia garantiza que la notación 
pueda cubrir nuevos escenarios técnicos. Si una notación deja de actualizarse, su utilidad académica y práctica 
disminuye, afectando directamente su adopción en el proceso formativo de los estudiantes. Por ello, la vigencia 
constituye un criterio clave en la toma de decisiones.
 

\subsubsection{Minimalismo}
\textbf{Código:} MINIMALISMO

\textbf{Descripción:} Este criterio evalúa qué tan independiente es la notación respecto a lenguajes e idiomas específicos,
así como el grado de simplicidad y neutralidad de sus símbolos. Se privilegia a las notaciones que permiten representar 
conceptos sin depender de terminología demasiado particular o propia de un idioma.


\textbf{Justificación:} Una notación minimalista favorece el aprendizaje y permite que los estudiantes comprendan los 
conceptos fundamentales sin verse limitados por implementaciones o el idioma. Para el área de infraestructura, 
esta independencia facilita que la notación pueda aplicarse en cursos diversos, proyectos y escenarios. Por estas 
razones, el minimalismo resulta un criterio apropiado al evaluar opciones dentro del proyecto NDITI.

\subsection{Análisis DAR}
Con los criterios definidos y la medición de popularidad completada, se procede a realizar el análisis \DAR. Para este análisis se utilizó una metodología estructurada basada en la plantilla encontrada en la tabla ~\ref{tab:plantilla-analisis_dar} que permite evaluar sistemáticamente cada alternativa contra los criterios establecidos. La plantilla utilizada se presenta en el Apéndice \ref{chap:plantilla-dar}.

Los resultados del análisis se presentan en la Tabla \ref{tab:analisis_notaciones}, donde cada notación fue evaluada con puntuaciones del 1 al 3 según su desempeño en cada criterio.

\begin{table}[H]
	\centering
	\renewcommand{\arraystretch}{1.2}
	\fontsize{9pt}{10pt}\selectfont
	\caption{Análisis de notaciones según criterios definidos.}
	\label{tab:analisis_notaciones}
	\begin{tabular}{|>{\centering\arraybackslash}p{2cm}
	                |>{\centering\arraybackslash}p{1.5cm}
	                |>{\centering\arraybackslash}p{1.8cm}
	                |>{\centering\arraybackslash}p{2.5cm}
	                |>{\centering\arraybackslash}p{1.4cm}
	                |>{\centering\arraybackslash}p{1.6cm}
	                |>{\centering\arraybackslash}p{1.4cm}|}
		\hline
		{\scriptsize\textbf{Notación}} &
		{\scriptsize\textbf{Popularidad}} &
		{\scriptsize\textbf{Documentación}} &
		{\scriptsize\textbf{Apalancamiento y Formalidad}} &
		{\scriptsize\textbf{Vigencia}} &
		{\scriptsize\textbf{Minimalismo}} &
		{\scriptsize\textbf{Promedio}} \\
		\hline
		
		\cellcolor[HTML]{E6E6E6}{\scriptsize\textbf{Conjunto}} &
		\multicolumn{5}{c|}{\cellcolor[HTML]{E6E6E6}{\scriptsize\textbf{Puntaje}}} &
		\cellcolor[HTML]{E6E6E6}{\scriptsize\textbf{ }} \\
		\hline
		
		BPMN & 3 & 2 & 3 & 1 & 1 & \cellcolor[HTML]{DBF7FF}2.0 \\
		\hline
		SysML & 1 & 2 & 3 & 2 & 1 & 1.8 \\
		\hline
		UML & 3 & 3 & 3 & 2 & 1 & \cellcolor[HTML]{FFCCC2}2.4 \\
		\hline
		ArchiMate & 1 & 3 & 1 & 3 & 2 & 2.0 \\
		\hline
		
	\end{tabular}
	\vspace{5pt}
\end{table}

\section{Resultados}

\subsection{Resumen del análisis}
El análisis realizado tuvo como propósito identificar qué notación ofrece un mejor equilibrio entre popularidad, disponibilidad de documentación, formalidad, vigencia y minimalismo. De acuerdo con los puntajes presentados en la Tabla \ref{tab:analisis_notaciones}, las notaciones con mejor desempeño relativo fueron \textbf{UML}, \textbf{BPMN} y \textbf{ArchiMate}, con promedios cercanos entre sí. Aunque UML obtuvo el puntaje más alto en términos generales, las diferencias frente a BPMN y ArchiMate no fueron lo suficientemente amplias como para descartar estas alternativas de manera concluyente.

Este resultado indica que no existe una única notación claramente dominante para el contexto del proyecto, sino un conjunto de opciones viables que aportan fortalezas distintas. UML destaca por su amplio respaldo documental y uso consolidado; BPMN por su orientación operativa y adopción extendida en modelado de procesos; y ArchiMate por su formalismo y vigencia en el ámbito de la arquitectura empresarial. En ausencia de una brecha significativa entre los puntajes, se opta por considerar las tres notaciones como candidatas válidas para las siguientes etapas del trabajo.

\subsection{Selección de la notación principal}
El análisis comparativo muestra que no existe una notación claramente dominante. UML obtuvo el mayor puntaje (2.4), destacándose por su madurez y uso extendido, mientras que BPMN sobresale en el modelado de procesos y ArchiMate en la representación de arquitectura empresarial.

Sin embargo, la selección de la notación principal no depende únicamente del puntaje, sino del nivel de abstracción requerido por el proyecto. UML y BPMN se enfocan en aspectos específicos del sistema, mientras que ArchiMate permite una visión integrada de las capas de negocio, aplicación y tecnología.

Por esta razón, se adopta ArchiMate como notación principal, debido a su mayor adecuación al enfoque de arquitectura empresarial del proyecto.